\scanpage{39}
\mspage{30}
% archive sheet number: 39
% inv: text | top | donc
 donc
% inv: display | top | deleted Mulomb equivalence
\[
\deleted{\unsure{\mathrm{Mulomb}|\mathrm{Muloms}}(X)\mathrel{|\circ|}\mathrm{Mulomb}(Y)\Longleftrightarrow}
\]
% inv: display | upper third | Mulomb disjointness equivalences
\[
\begin{aligned}
\mathrm{Mulomb}(X)\mathrel{|\circ|}\mathrm{Mulomb}(Y)&\Longleftrightarrow\deleted{\unsure{M_{\overset{\circ}{\ll}X}}}\mathrm{Omb}^\circ(X)\cap\mathrm{Omb}^\circ(Y)=\varnothing\\
&\Longleftrightarrow\nexists Z\in M,\text{ avec }Z\overset{\circ}{\ll}X,\ Z\overset{\circ}{\ll}Y.
\end{aligned}
\]
% inv: deletion | upper third | i.e. ceci
\deleted{i.e. ceci :}

% inv: deletion | upper middle | Est-ce
\deleted{Est-ce}\deleted{\unsure{que}}
% inv: text | upper middle | On a donc bien
On a donc bien
% inv: display | middle | X disjoint Y implies
\[
X\mathrel{|\circ|}Y\Longrightarrow\mathrm{Mulomb}(X)\mathrel{|\circ|}\mathrm{Mulomb}(Y)
\]
% inv: insertion | below implication | i.e. X Y pas minoré
\inserted{i.e. $\{X,Y\}$ pas minoré \unsure{ds|dans} $M,\overset{\circ}{\ll}$.}

% inv: text | middle | Car si on avait
Car si on avait $Z\overset{\circ}{\ll}X$, $Z\overset{\circ}{\ll}Y$, on aurait $Z\deleted{\overset{\circ}{\ll}}\inserted{\mathrel{|\circ|}}Z$ par Mag 3 b), absurde par Mag 3 a).

% inv: text | middle lower | Mais l'inverse
Mais l'inverse est-il tjrs vrai, i.e. $X\mathrel{|\circ|}Y$ peut-il s'exprimer en termes de la relation d'ordre $\overset{\circ}{\ll}$ ?

% inv: deletion | lower middle | Ça ne marche que
\deleted{Ça \unsure{ne marche que}} \deleted{C'est OK dans}

% inv: text | lower third | Mais notons que
Mais notons que si $x,y\in L$, $x\ne y$, alors $\{x,y\}$ n'est pas minoré dans $M$ pour $\overset{\circ}{\ll}$, mais on n'aura pas pour autant, en général, $x\mathrel{|\circ|}y$\deleted{\unsure{?}} i.e. $x\parallel y$ ! (Même, dans le cas des magasins \unsure{dérivables|divisibles}.)

% inv: text | bottom | Donc le fait
Donc le fait que $\mathrel{|\circ|}$ soit induit par $\overset{\circ}{\ll}$ s'apparente à une hypothèse de nature ``quasi-ensembliste'', mais sans référence particulière aux lieux.
